% Cadena diagnóstica de cuatro comprobaciones ordenadas.
% Cada comprobación que falla identifica una causa distinta.
\begin{singlespace}
\begin{tikzpicture}[figbase]

    % ------------------------------------------------- columna de comprobación
    \node[figinfra, minimum width=4.2cm, minimum height=0.95cm] (ini) at (0,0)
        {Intent evaluado sobre un perfil};

    \node[figgate] (g1) at (0,-2.05)
        {\textbf{Materialización}\\[-1pt]{\scriptsize ¿el discriminador ya existe?}};
    \node[figgate] (g2) at (0,-4.35)
        {\textbf{Observación confiable}\\[-1pt]{\scriptsize ¿con procedencia suficiente?}};
    \node[figgate] (g3) at (0,-6.65)
        {\textbf{Expresividad}\\[-1pt]{\scriptsize ¿el lenguaje lo representa?}};
    \node[figgate] (g4) at (0,-8.95)
        {\textbf{Acción}\\[-1pt]{\scriptsize ¿actúa antes del productor?}};

    \node[figpos] (r) at (0,-11.25)
        {\textbf{R}\\[-1pt]{\scriptsize realizable}};

    % Distintivo numerado sobre el borde de cada comprobación.
    \foreach \i/\g in {1/g1, 2/g2, 3/g3, 4/g4}
        \node[circle, draw=figNFborde, fill=white, text=figNFtexto,
              line width=0.7pt, font=\scriptsize\bfseries,
              inner sep=0pt, minimum size=4.6mm] at (\g.west) {\i};

    % ------------------------------------------------------- causas de fracaso
    \node[figneg] (f1) at (7.1,-2.05)
        {\textbf{U (F1)} \enspace La información requerida todavía no existe en el
         punto de decisión.};
    \node[figneg] (f2) at (7.1,-4.35)
        {\textbf{U (F2)} \enspace Existe, pero el PEP no la observa con la confianza
         necesaria: no es visible, la oculta el cifrado o la escribe un actor no
         confiable.};
    \node[figneg] (f3) at (7.1,-7.80)
        {\textbf{U (F3)} \enspace La información confiable está disponible, pero el
         lenguaje no representa el predicado o el mecanismo no ejecuta la acción
         preventiva.};

    % ------------------------------------------------------------------ flujos
    \draw[figflujo] (ini) -- (g1);
    \foreach \a/\b in {g1/g2, g2/g3, g3/g4}
        \draw[figsi] (\a) -- (\b) node[figrot, midway, right=2pt, text=figPosTexto] {sí};
    \draw[figsi, line width=1.2pt] (g4) -- (r)
        node[figrot, midway, right=2pt, text=figPosTexto] {sí};

    \draw[figno] (g1.east) -- (f1.west) node[figrot, midway, above=1pt, text=figNegTexto] {no};
    \draw[figno] (g2.east) -- (f2.west) node[figrot, midway, above=1pt, text=figNegTexto] {no};
    \draw[figno] (g3.east) -- ++(1.15,0) |- ([yshift=0.55cm]f3.west);
    \draw[figno] (g4.east) -- ++(1.15,0) |- ([yshift=-0.55cm]f3.west);
    \node[figrot, text=figNegTexto] at ($(g3.east)+(0.52,0.18)$) {no};
    \node[figrot, text=figNegTexto] at ($(g4.east)+(0.52,0.18)$) {no};

    % --------------------------------------------- comprobación independiente
    \node[fignota, text width=13.3cm] (nota) at (4.0,-13.35)
        {\textbf{Realización en ejecución.}\enspace Una petición permitida puede fallar
         después por indisponibilidad del productor o por un problema de ruta. Ese
         resultado se registra como \texttt{RUNTIME\_REALIZATION\_FAILURE} y no
         modifica el veredicto anterior.};

    \draw[figaux] (r.south) -- (r.south |- nota.north);

\end{tikzpicture}
\end{singlespace}
